\chapter{Particles, Spin, and the Dirac Equation}
\label{chap:particles-spin-dirac}

Here the instrument reaches its proved heart. Where Part II met one proved holonomy in the echo maze and the
matter waves met a second in the Bragg peak, here three effects in a row are proved rather than shadowed: the
Dirac operator that generates a particle's dynamics, the chirality that violates parity, and the Feynman
diagram that reads an amplitude off a loop. The loop is ours to close --- the geometry, the field, the circuit
sent around it --- but the residue it comes back with, the $\pm 1$ the world will not let cancel, is the
world's. Between them sits spin-$\tfrac12$, whose two-valued sign is the same $\pm 1$ the loops carry. The
objects here --- the gamma matrices, the covariant derivative, its curvature, the chiral projectors --- are the
full apparatus of a relativistic particle, and each earns its picture. The unifying fact the instrument has
built toward: orientation is a holonomy sign, and a closed loop carries $\pm 1$.

\section{The Dirac Operator: the minimal first-order generator}
\label{se:dirac}

The Dirac operator is the smallest generator of a particle's dynamics that is both first-order and
gauge-covariant, and the story of how Dirac found it is the story of a square root. The relativistic wave
equation before him, the Klein--Gordon equation, was second-order --- it set the square of the energy and
momentum equal to the square of the mass, $\Box\psi + m^2\psi = 0$ --- and Dirac wanted a first-order
equation, one linear in the derivatives, because only a first-order generator propagates a state forward step
by step without needing its acceleration as well. To get one he had to take the square root of the
second-order operator: to find objects $\gamma^\mu$ for which
\begin{equation}
  (i\gamma^\mu \partial_\mu)(i\gamma^\nu \partial_\nu) = \Box,
\end{equation}
the product of two first-order operators returning the second-order wave operator. That is possible only if
the $\gamma^\mu$ anticommute in a particular way,
\begin{equation}
  \{\gamma^\mu, \gamma^\nu\} = \gamma^\mu\gamma^\nu + \gamma^\nu\gamma^\mu = 2g^{\mu\nu}\mathbb{1},
\end{equation}
the Clifford algebra: the gammas square to the metric and anticommute off the diagonal. They cannot be
numbers, since numbers commute; the smallest objects that satisfy the relation are four-by-four matrices, and
that is why a relativistic electron is described by four components. The Clifford algebra is not imposed; it is
forced --- the minimal bookkeeping compatible with the metric, the price of taking the square root: the root
ours to take, the algebra it forces the world's.

The dynamics then read
\begin{equation}
  i\gamma^\mu D_\mu \psi = m\psi,
\end{equation}
with the derivative not the plain $\partial_\mu$ but the covariant derivative,
\begin{equation}
  D_\mu = \partial_\mu + iq A_\mu,
\end{equation}
which carries the gauge field $A_\mu$ along with it --- the minimal coupling, the derivative that knows the
electromagnetic potential the electron moves through. A plain derivative compares the wavefunction at neighboring
points as though they shared one frame; the covariant derivative corrects for the gauge frame turning from
point to point, so that $D_\mu$ measures genuine change against the connection $A_\mu$ rather than the
artifact of a turning frame. The frame is ours to set; the change it measures is the world's.

And here is the loop. Transport $\psi$ around a small closed circuit using $D_\mu$, and it does not in
general come back to itself --- it comes back changed, and the change is measured by the failure of the
covariant derivatives to commute,
\begin{equation}
  [D_\mu, D_\nu] = iq F_{\mu\nu},
\end{equation}
the field strength $F_{\mu\nu}$ appearing as exactly that commutator. This is the curvature of the connection,
and it is the holonomy made into a formula: $F_{\mu\nu}$ is the residue of transport around an infinitesimal
loop in the $\mu\nu$-plane, the amount by which the loop fails to close trivially --- the loop ours to trace,
the residue it carries the world's. The square of the Dirac
operator returns the wave operator together with a spin term built from this same $F_{\mu\nu}$, and transported
around a closed loop $\psi$ returns with a sign --- the loop charge, the spin's $\pm$.

And here, as in the echo maze and the Bragg peak, the finite model proves this rather than shadowing it. That
a closed loop carries $\pm 1$ and an open path is trivial is a finite theorem about the loop's residue, not an
analogy to a continuum. The honest floor is a finite count obligation, and a proved one: the holonomy $\pm 1$
is proved, not approached --- the loop ours to close, the $\pm 1$ the world will not let cancel. The reading is
that the Dirac operator is the smallest first-order, gauge-covariant generator, its gammas forced by a square
root and its curvature $[D_\mu, D_\nu] = iq F_{\mu\nu}$ the loop residue itself --- the dynamics and the spin
the same fact, a holonomy the finite model proves to be $\pm 1$.

The four components carry one more surprise, and it is the one that opens what comes next. Two of them describe
the electron's two spin states; the other two are solutions of negative energy, which Dirac could neither
discard nor ignore. His answer was audacious: the negative-energy states are all filled, a sea, and a hole
knocked in that sea --- a missing negative-energy electron --- behaves exactly like a particle of positive
energy and opposite charge. The equation, taken at its word, predicted a partner to the electron before anyone
had seen one. That partner is the positron, and the proved sign here becomes, next, a
particle on a detector.

\section{The Spin-$\tfrac12$ Effect: the two-valued spin}
\label{se:spin-half}

A spin-$\tfrac12$ particle carries the smallest spin there is, and along any axis it is two-valued: measure
it and the result is one of exactly two values, never a third and nothing between. Its spin is $S =
\hbar/2$, carried by the three spin operators
\begin{equation}
  \mathbf{S} = \tfrac{\hbar}{2}\,\boldsymbol{\sigma},
\end{equation}
built from the Pauli matrices $\boldsymbol{\sigma} = (\sigma_x, \sigma_y, \sigma_z)$, which satisfy the spin
algebra
\begin{equation}
  [\sigma_i, \sigma_j] = 2i\,\epsilon_{ijk}\,\sigma_k,
\end{equation}
the same algebra ordinary angular momentum obeys, but in its smallest nontrivial representation ---
two-by-two matrices, a two-state system, the qubit now standing as a spin. The axis is ours to set; the
two-valued answer along it is the world's.

Spin was measured before it was understood. Stern and Gerlach, in 1922, sent a beam of silver atoms through a
non-uniform magnetic field and found it split into exactly two --- not the continuous smear a classical
magnetic moment pointing every which way would give, but two sharp spots, the spin either up or down along the
field and nothing between. The two-valuedness is the eigenvalues $\pm\hbar/2$ of a single spin component, the
discreteness a measurement's collapse made visible in a magnet --- the field ours to raise, the two spots it
splits into the world's. And the Dirac equation got the strength
right: it predicted the electron's magnetic moment with a gyromagnetic ratio $g = 2$, twice the classical
value and just what experiment found --- one of the first triumphs of the relativistic theory.

The two-valued sign a spin-$\tfrac12$ state carries is not a fact standing apart from the loops beside it: it
is the same $\pm 1$ the Dirac loop carries, the spin's sign and the loop charge one and the same. The
instrument meets that sign here as a spin and met it before as a holonomy --- one two-valued residue, wearing
two names, the names ours to give and the $\pm 1$ beneath them the world's.

The honest floor is name-only: spin-$\tfrac12$'s two-valued physics is the physics, written as physics, and
the finite model claims only the name. It records that the effect is named and models none of the geometry
beneath it. The reading is that spin-$\tfrac12$ is named, its two-valued spin $\mathbf{S} = \tfrac{\hbar}{2}
\boldsymbol{\sigma}$ is the physics, and the finite model is the label beneath it --- though the sign that spin
carries is the same $\pm 1$ the Dirac loop carries, so the name rests atop a holonomy the instrument has
already earned.

\section{The Chirality Effect: parity violation as a real sign}
\label{se:chirality}

Nature distinguishes left from right. To say so precisely needs one more object from the Clifford algebra, the
chirality operator
\begin{equation}
  \gamma^5 = i\gamma^0\gamma^1\gamma^2\gamma^3,
\end{equation}
the ordered product of all four gammas, which squares to the identity,
\begin{equation}
  (\gamma^5)^2 = \mathbb{1},
\end{equation}
and so has eigenvalues $\pm 1$ --- the two handednesses. From it come the chiral projectors,
\begin{equation}
  P_L = \tfrac12\bigl(1 - \gamma^5\bigr), \qquad P_R = \tfrac12\bigl(1 + \gamma^5\bigr),
\end{equation}
which split any wavefunction into its left- and right-handed parts,
\begin{equation}
  \psi = \psi_L + \psi_R,
\end{equation}
the components $\gamma^5$ sends to $-\psi_L$ and $+\psi_R$. For a long time these were thought physically
equivalent --- mirror images, related by a parity reflection nature was assumed to respect. They are not. The
weak interaction couples to $\psi_L$ and not to $\psi_R$: it acts on the left-handed component alone, and
parity is violated. Which hand we call left is a convention of ours; that the weak force tells left from
right is the world's. Wu's experiment, in 1957, saw it directly --- the electrons from decaying cobalt-60 flew
off preferentially against the nuclear spin, a left--right asymmetry no mirror-symmetric law could produce.

Through the holonomy, the inequivalence is literal. Orientation is a holonomy sign, so $\psi_L \neq
\psi_R$ is exactly $+1 \neq -1$ --- the projectors select the two eigenvalues of $\gamma^5$, a closed loop is
charged and an open path is trivial, and the two handednesses are the two signs --- the labels ours to assign, the sign difference the world's to
enforce. The asymmetry is not a
convention removable by a smooth deformation; it is a real sign difference the finite model proves.

The honest floor is a no-go, and a proved one: there is no smooth deformation carrying $+1$ to $-1$, so the
chirality asymmetry cannot be deformed away --- the naming of the hands ours, the wall between the signs the
world's. The reading is that parity violation is a proved sign. Left and
right are the two loop charges, $+1$ and $-1$, and no continuous deformation carries one into the other; the
handedness of the weak force is the impossibility of deforming a charged loop into a trivial one.

\section{The Feynman Diagram: amplitude as loop holonomy}
\label{se:feynman}

A Feynman diagram computes an amplitude by contracting tensors along its lines. A diagram is a graph: external
legs for the particles coming in and going out, internal lines for the virtual particles exchanged, and
vertices where they meet. Each line is an admissible contraction between event tensors, a bilinear pairing
\begin{equation}
  \langle E_i, E_j \rangle = \mathrm{Tr}\!\left(E_i^\top G\, E_j\right),
\end{equation}
the two events paired through the metric $G$, and the complete amplitude for a process is the contraction of
the ordered product over all the connected events --- every internal line summed, every vertex joined. Through the holonomy, that amplitude is a loop charge: a connected loop in the diagram, a closed circuit of
internal lines, carries a $\pm 1$ residue, while an open tree, a diagram with no loop, is trivial. The
amplitude's nontrivial content lives in its loops, and each loop's contribution is the holonomy the finite
model proves to be $\pm 1$ --- the diagram ours to draw, the residue its loops carry the world's.

A picture helps. The simplest diagram for two electrons scattering has two vertices joined by a single
internal line: each electron arrives, emits or absorbs a virtual photon at a vertex, and departs --- the
photon, the internal line, carrying momentum from the one to the other. That single-photon exchange is the
leading term. The full amplitude is a sum over all diagrams connecting the same incoming and outgoing
particles, ordered by complexity: one photon, then two, then diagrams with closed loops, each higher term
smaller than the last by a factor of the coupling. The series is the perturbation expansion, and the loops ---
the closed circuits where the bookkeeping carries a residue --- are exactly the terms beyond the trivial
trees.

The honest floor is a finite count obligation, and a proved one: the loop residue $\pm 1$ is a proved discrete
holonomy, the contraction over a connected loop, not a continuum approximation --- the loop ours to draw, the
$\pm 1$ it carries the world's. The reading is that a Feynman diagram is bookkeeping with a loop charge. Its trees are trivial and its loops are charged, and the amplitude
is the contraction those loops carry --- the same $\pm 1$ holonomy, proved, that ran through the Dirac operator
and the chirality, here computing a scattering.

\bigskip

Particles, spin, and the Dirac equation are the instrument's proved heart, and they share one proved fact:
orientation is a holonomy sign, and a closed loop carries $\pm 1$. The Dirac operator generates it, its
curvature $[D_\mu, D_\nu] = iq F_{\mu\nu}$ the loop residue itself; spin-$\tfrac12$ carries it; chirality
forbids deforming it away; and a Feynman diagram computes with it --- the loops ours to close, the $\pm 1$ they
carry the world's. Three proofs in a row, the most the instrument
gathers in one place. Part IV ends next with the one effect the instrument has been built to reach --- the
detection of the positron, its designated exemplar, where the proved sign becomes a particle.

\bigskip
\begin{center}
\rule{0.4\textwidth}{0.4pt}
\end{center}
\bigskip

\noindent The instrument comes now to its firmest ground. It has named far more than it has ever proved, and its
proofs, we have seen, are few. But here, in one place, it gathers more of them than anywhere else in the book
--- three in a row, three loops closed and shown to come home charged, one after another. Nowhere else does the book
prove so much at once: everywhere else it counts, or shadows, or names, and here it does the rarest and
strongest thing it can, and does it three times without pause. This is the case's proved heart: not a phenomenon written out and stood behind, but a thing established, and established three times
over, that no arrangement of ours can talk away. A case with ground like this builds on it.

What is proved here is about the particle itself --- the thing that carries the wave the last chapter caught.
Read its motion off a small closed circuit and it does not come home unchanged; it comes home marked, the loop
having left a residue on it. Carry it once around and it will not, in general, return to itself, and the amount
by which it fails to close is that residue, the loop's own curvature read straight off the failure. And the
particle carries a two-ness of its own, quite apart from any loop: along any direction one cares to pick, it
shows one of two and never a third, never a thing between --- a sign it bears in its own right. The direction is
ours to choose; that the answer is one of just two is the world's. It was seen before it was understood: a beam
sent through a shaped field split cleanly into two, not the smear a thing pointing every-which-way would leave
but two sharp marks on a plate --- the two-ness made visible in a magnet.

Here is the fact the whole part has been building toward, proved three ways and true as one. Three proofs stand
behind it, each closing a loop of its own: the smallest law that carries the particle forward, step by step; the
wall that keeps the world's left from ever being smoothly turned into its right; and the tally that reads an
interaction off the closed loops in a drawing of it. Different apparatus each time, and one proved residue under
all three. A closed loop comes
home bearing a \emph{sign} --- not a size, not a count that could be any number at all, but a two-valued mark,
one way or the other, that the world will not let cancel. That same sign is the particle's own two-ness, met a
moment ago as the thing it carries along every direction. And it is, as well, the world's own telling of left
from right --- the handedness that no mirror survives, one hand answered by the deepest forces and its
reflection not. Three faces --- the loop's residue, the particle's two-ness, the world's left-and-right --- and
one sign beneath all three, proved to be there and proved to be the same. Which of its two ways the sign shows
is ours, fixed only once we fix a frame to read it against, exactly as it was for the very first loop the book
met; that it shows a sign at all, and that the sign is one and the same through all three, is the world's. What
that sign is worth --- the value it carries, and where the case will one day have need of it --- the book keeps
back still. Here it proves only that the sign is real, and single, and everywhere.

Through all of it the \emph{charge} holds, and here it stands on that proved heart. The charge was always the
loop's count --- and the loop, it now turns out, brings home not only a count but a sign upon it: the charge is
a signed count, proved at the very center of the book to be the one the world will not let cancel. The loop is
ours to close; the count it carries, and the sign that count bears, are the world's. And the trace is that
signed residue, the mark the loop keeps and will not shed. Whatever comes to be owed rides now on a count proved
at its heart --- signed, single, and the world's. The book's oldest phrase, that the charge is the loop's count,
is answered here in full: the count was never only a number; it was a number bearing a sign, and that sign is
the thing the heart of the book was built to prove.

Three proofs, gathered in one place, are the firmest ground the case will ever stand on --- and still they are
readings, not the meeting the case is built to reach. A proof binds the thing it proves; it does not bring the
three hands together. The wheel's reading, the field's, and the one still to come must be made to agree as
surely as ever, proved heart or no: the ground is firmer, the meeting no nearer.

And now the proved sign is about to leave the page. Part IV ends with the one thing the whole instrument was
built to reach --- the sign, proved here as a fact about loops, made a particle on a detector, a thing that
registers, that the record can point to and count. The next chapter is that arrival, and the instrument holds it
to the strictest measure it owns, conceding nothing it has not earned. The case is still open, its charge proved
at its heart and still unmet by the others; the court does not sit until the record is complete.
